% ECONOMICS_PROBLEM: {"id": "C73-606", "title": "Learning to sustain cooperation", "jel_code": "C73", "source_id": "OP-0282"}
% !TeX program = xelatex
\documentclass[11pt,a4paper]{article}
\usepackage[margin=23mm]{geometry}
\usepackage{fontspec}
\setmainfont{Latin Modern Roman}
\usepackage{amsmath,amssymb,mathtools}
\usepackage{xcolor,enumitem,booktabs,longtable,array,xurl,hyperref}
\definecolor{muted}{HTML}{53636C}
\hypersetup{colorlinks=true,urlcolor=blue,linkcolor=blue}
\setlength{\parindent}{0pt}
\setlength{\parskip}{5pt}
\newcommand{\R}{\mathbb R}
\newcommand{\E}{\mathbb E}
\newcommand{\N}{\mathbb N}
\newcommand{\ind}{\mathbf 1}
\newcommand{\Var}{\operatorname{Var}}
\newcommand{\Cov}{\operatorname{Cov}}
\newcommand{\supp}{\operatorname{supp}}
\DeclareMathOperator*{\argmin}{arg\,min}
\DeclareMathOperator*{\argmax}{arg\,max}
\newcommand{\entrymeta}[3]{{\sffamily\small\color{muted}\textbf{\texttt{#1}}\quad|\quad #2\par #3\par}}
\newcommand{\Problemref}[1]{\texttt{#1}}
\newcommand{\JELref}[2]{\texttt{#2} (JEL \texttt{#1})}
\begin{document}
\section*{Learning to sustain cooperation}
\textbf{Problem ID:} \texttt{C73-606}\quad \textbf{JEL:} \texttt{C73}\par
% BEGIN ECONOMICS STATEMENT
\label{problem:OP-0282}
{\sffamily\small\color{muted}Original subject group: Game theory and strategic interaction\par}
\label{definitions:problem:30007564}
\textbf{Audit classification:} Published research agenda. Fudenberg–Levine pp.162–165 discusses off-path learning and cooperation. The experimental question is editorial; the original formula double-counted random continuation.
\subsection*{Definitions and precise research target}
\textbf{Model and research target.} In the explicitly specified repeated prisoner’s dilemma, determine how experimentally provided information about off-path punishment and forgiveness changes the correctly normalized cooperation measure \(C(1)-C(0)\). Estimate its dependence on continuation \(\beta\), the full monitoring kernel and elicited beliefs; fit a prespecified structural learning model and evaluate its forecasts in new monitoring regimes. This is an empirical research specification, not a universal equilibrium-existence conjecture.

There are two players, each choosing \(C\) or \(D\) every surviving period. Stage payoffs are \((2,2)\) at \((C,C)\), \((1,1)\) at \((D,D)\), \((0,3)\) at \((C,D)\), and \((3,0)\) at \((D,C)\). \emph{Definitions and maintained assumptions (editorial unless a source definition is named).} Let \(N\) be the last played period, with \(\Pr(N\ge t)=\beta^{t-1}\), and suppose termination draws are independent of actions and monitoring. Define \(Y_t=\mathbf1\{a_{1t}=a_{2t}=C\}\) on surviving histories and use
\[C(Z)=(1-\beta)\mathbb E\!\left[\sum_{t=1}^{N}Y_t\mid Z\right]=(1-\beta)\sum_{t\ge1}\beta^{t-1}\mathbb E[Y_t\mid N\ge t,Z].\]
Do not multiply observed surviving outcomes by another \(\beta^{t-1}\): survival already supplies that factor. Then \(0\le C(Z)\le1\). The private monitoring kernel must specify each player’s signal conditional on the complete action profile; marginal error probability \(\rho\) alone does not specify joint errors. For a concrete baseline take independent action-misclassification errors conditional on actions. Treatment changes only the declared feedback kernel and is assigned by pair. A learning model specifies beliefs over opponents’ contingent strategies and update and choice rules at every private history. Incentivized belief reports require a scoring rule and error model.
\subsection*{Variants and assumption changes}
\begin{enumerate}
\item \textbf{Public versus private monitoring (editorial).} Randomize a common public action signal versus independent private signals with matched marginal accuracy. Estimate the cooperation contrast and test whether a fitted off-path learning model predicts it.
\item \textbf{Pure time discount (editorial).} Add a discount factor \(\delta\) distinct from survival. Use \((1-\beta\delta)\sum_t(\beta\delta)^{t-1}\mathbb E[Y_t\mid N\ge t,Z]\), and specify the corresponding estimator for terminated histories.
\end{enumerate}
\subsection*{References and source audit}
\begin{enumerate}
\item Off-Path Play pp.162–163; Cooperation pp.163–165. Supports the stated research area; exact displayed specification is editorial. \url{https://economics.mit.edu/sites/default/files/2022-09/Whither%20Game%20Theory%20Towards%20a%20Theory%20of%20Learning%20in%20Games.pdf}
\end{enumerate}
\begin{enumerate}\raggedright
\item\hypertarget{definitions:source:30007564:1}{} Drew Fudenberg; David K. Levine. 2016. \emph{Whither Game Theory? Towards a Theory of Learning in Games}. Active Learning: Learning about Off-Path Play, pp.162–163; Cooperation in Infinitely Repeated Games, pp.163–165.
\url{https://economics.mit.edu/sites/default/files/2022-09/Whither%20Game%20Theory%20Towards%20a%20Theory%20of%20Learning%20in%20Games.pdf}.
DOI: \url{https://doi.org/10.1257/jep.30.4.151}.
\end{enumerate}

\subsection*{Common conventions: Source mathematical conventions}

These conventions apply to each entry unless explicitly replaced.

\textbf{Models and identification.} A primitive model \(m\) specifies all probability laws, feasible choices, information, equilibrium and measurement rules used by its target. The admissible class \(\mathcal M\) is fixed independently of outcomes. If \(P_O(m)\) is its observable probability law and \(\tau(m)\) the target, its identified set at observable law \(P\) is
\[
 \mathcal I_\tau(P)=\{\tau(m):m\in\mathcal M,\ P_O(m)=P\}.
\]
Point identification means that this set is a singleton. An empty set signals incompatibility of data and assumptions. Coordinate bounds need not describe the joint set. Bounds are sharp only if every retained target is attainable or arbitrarily approachable by compatible models. Finite bounds require explicit restrictions on unobserved quantities; unrestricted confounding, hidden wealth, extrapolation or arbitrary equilibrium selection can yield uninformative sets.

\textbf{Causal interventions.} A potential outcome \(Y_i(p)\) is the outcome under a complete policy regime \(p\), including timing, financing, information and market spillovers. The target population and units are fixed before treatment. With interference, use the complete assignment vector or a justified exposure mapping; individual no-interference notation alone cannot identify an economy-wide intervention. A difference of mean potential outcomes does not require the cross-world joint distribution. Conditioning on post-treatment migration, survival, enrollment or employment changes the estimand and needs separate justification.

\textbf{Statistical statements.} An estimator, confidence set, forecast or test specifies a sampling law, model class and loss/coverage criterion. Uniform \((1-\alpha)\)-coverage means \(\inf_{m\in\mathcal M}\Pr_m\{\tau(m)\in C_n\}\geq1-\alpha\), or an explicitly asymptotic version. The whole parameter space trivially has coverage, so informativeness requires a stated width, risk or power objective. A forecasting improvement is relative to a named benchmark, information set and evaluation population.

\textbf{Equilibrium and optimization.} An equilibrium comprises feasible plans, optimality, beliefs where needed, budget constraints and all market/resource clearing. The primitives or a stated benchmark define each correspondence \(\mathcal E(p;m)\). Nonemptiness is assumed or proved, never inferred just from compact policy sets. A selected-equilibrium target uses a declared selection rule; a robust target quantifies over all permitted equilibria. Use suprema or infima unless attainment is established. An \(\varepsilon\)-optimal policy is feasible and within \(\varepsilon\) of the finite optimum value. Report an empty feasible set separately.

\textbf{Welfare and accounting.} Normative utility, interpersonal normalization, social weights, numeraire, discounting and the use of public receipts are primitives. Behavioral utility need not coincide with the welfare criterion. Household welfare includes ownership income and rents once; adding profits again to compensating variation generally double-counts them. A monetary equivalent is defined by an explicit transfer and price convention, with monotonicity and range conditions. Average effects, mechanism contributions, transfers and welfare losses are different targets. Infinite sums require convergence and dynamic feasible plans require terminal or no-Ponzi conditions.

\textbf{Source versions.} A working-paper date, revision date and journal publication year are distinguished. A DOI for a journal version is not automatically the DOI of an earlier manuscript. Locators refer to the stated version and printed pages where specified. A metadata-only check does not verify a claim about a full-text conclusion. Every entry's source subsection narrows the provenance claim accordingly.
% END ECONOMICS STATEMENT
\end{document}
